\section{Результаты моделирования}

\subsection{Параметры модели}

Моделирование проведено для радиального подшипника скольжения геометрии~\varGeomLetter{} ($R = \varRmm$~мм, $c = \varCmkm$~мкм, $L = \varLmm$~мм) с \varDimpleTypeGenitive{} микрорельефом. Исходные данные сведены в таблицу~\ref{tab:input_data}.

\begin{table}[H]
\small
\caption{Исходные данные для расчёта}
\label{tab:input_data}
\begin{tabularx}{\textwidth}{|X|c|c|}
\hline
\centering\textbf{Параметр} & \textbf{Значение} & \textbf{Размерность} \\
\hline
Радиус вала $R$ & \varRmm & мм \\
\hline
Радиальный зазор $c$ & \varCmkm & мкм \\
\hline
Длина подшипника $L$ & \varLmm & мм \\
\hline
Относительная длина $L/D$ & \varLD & -- \\
\hline
Полуось по оси $Z$, $a$ & \varDimpleAmm & мм \\
\hline
Полуось по оси $\varphi$, $b$ & \varDimpleBmm & мм \\
\hline
Глубина углубления $h_p$ & \varDimpleHpmkm & мкм \\
\hline
Относительная глубина $h_p/c$ & \varDimpleHpRatio & -- \\
\hline
Число углублений по окружности $N_\varphi$ & \varNphi & шт. \\
\hline
Число рядов по длине $N_Z$ & \varNZ & шт. \\
\hline
Область текстурирования & \varPhiRange & -- \\
\hline
Частота вращения $n$ & \varN & об/мин \\
\hline
Вязкость смазки $\eta$ & \varEta & Па$\cdot$с \\
\hline
Диапазон эксцентриситета $\varepsilon$ & 0{,}05--0{,}80 & -- \\
\hline
Размер сетки & $500 \times 500$ & -- \\
\hline
Параметр релаксации $\omega$ & 1{,}5 & -- \\
\hline
Критерий сходимости $\delta$ & $< 10^{-5}$ & -- \\
\hline
Макс. число итераций & 50000 & -- \\
\hline
\end{tabularx}
\end{table}

Линейная скорость поверхности вала при заданной частоте вращения равна $U = \pi D n / 60 = \pi \cdot 0{,}100 \cdot 2980 / 60 \approx 15{,}60$~м/с. Расчёт проведён для 15 значений эксцентриситета в указанном диапазоне.

\subsection{Распределение зазора и давления}

Распределение безразмерного давления $P(\varphi)$ в среднем сечении ($Z = 0$) при $\varepsilon = 0{,}6$ для гладкого и текстурированного подшипников показано на рисунке~\ref{fig:pressure_2d}.

\begin{figure}[H]
\centering
\includegraphics[width=\textwidth]{figures/pressure_2d.png}
\caption{Распределение давления $P(\varphi)$ при $Z=0$ для гладкого и текстурированного подшипников ($\varepsilon = 0{,}6$)}
\label{fig:pressure_2d}
\end{figure}

В гладком подшипнике давление имеет один максимум в зоне минимального зазора ($\varphi \approx \pi$). Оно возрастает в сужающейся части зазора (от $\varphi = 0$ до $\varphi \approx \pi$) и обращается в нуль в расширяющейся части из-за кавитации. Зона положительного давления охватывает приблизительно половину окружности. Текстурирование суперэллиптическими углублениями вносит локальные возмущения в кривую давления: на входной кромке каждого углубления формируется пик давления, а на выходной~--- провал. Суперэллиптический профиль, описываемый степенной функцией четвёртого порядка $\Delta h = h_p \sqrt{1 - r^4}$, обладает более крутыми стенками по сравнению с эллипсоидальным (второго порядка), что создаёт более резкие перепады давления на кромках. Тем не менее общий уровень давления при наличии микрорельефа выше, что обеспечивает прирост несущей способности.

На рисунке~\ref{fig:clearance_2d} показан профиль зазора $H(\varphi)$ в среднем сечении при $\varepsilon = 0{,}6$.

\begin{figure}[H]
\centering
\includegraphics[width=\textwidth]{figures/clearance_2d.png}
\caption{Распределение зазора $H(\varphi)$ при $Z=0$ для подшипника с \varDimpleTypeGenitive{} углублениями ($\varepsilon = 0{,}6$)}
\label{fig:clearance_2d}
\end{figure}

На фоне плавного изменения зазора по закону $H = 1 + \varepsilon\cos\varphi$ видны локальные увеличения, соответствующие суперэллиптическим углублениям. По форме суперэллиптический профиль занимает промежуточное положение между эллипсоидальным и цилиндрическим: стенки углублений круче, чем у эллипсоида, но без резкого перехода, характерного для цилиндра. Глубина каждого углубления составляет $\Delta H_{\max} = h_p/c = 0{,}1$ в безразмерных единицах.

Трёхмерные распределения зазора и давления при $\varepsilon = 0{,}6$ представлены на рисунках~\ref{fig:clearance_3d} и~\ref{fig:pressure_3d}.

\begin{figure}[H]
\centering
\includegraphics[width=\textwidth]{figures/clearance_3d.png}
\caption{Трёхмерное распределение зазора $H(\varphi, Z)$ ($\varepsilon = 0{,}6$)}
\label{fig:clearance_3d}
\end{figure}

На рисунке~\ref{fig:clearance_3d} чётко прослеживается регулярная структура микрорельефа: 8 углублений по окружности и 11 рядов по длине подшипника в зоне $90^\circ - 270^\circ$.

\begin{figure}[H]
\centering
\includegraphics[width=\textwidth]{figures/pressure_3d.png}
\caption{Трёхмерное распределение давления $P(\varphi, Z)$ ($\varepsilon = 0{,}6$)}
\label{fig:pressure_3d}
\end{figure}

На рисунке~\ref{fig:pressure_3d} помимо основного максимума давления, обусловленного макрогеометрией клинового зазора, наблюдаются локальные возмущения давления вблизи каждого углубления. На входной кромке углубления (по направлению вращения) зазор резко сужается, что порождает дополнительный подъём давления~--- микрогидродинамический эффект. Благодаря более крутым стенкам суперэллиптического профиля (степень~4 вместо~2) перепад давления на кромках выражен сильнее, чем для эллипсоидальной формы. На выходной кромке давление снижается, однако кавитационное ограничение ($P \geq 0$) обеспечивает положительный суммарный вклад каждого углубления в несущую способность.

Анализ трёхмерных распределений позволяет выделить следующие особенности. Давление распределено неравномерно по длине подшипника~--- максимум приходится на среднее сечение ($Z = 0$), а к торцам ($Z = \pm 1$) давление спадает до нуля вследствие торцевых утечек. Влияние микрорельефа наиболее существенно вблизи среднего сечения, где базовое давление максимально. На периферийных рядах углублений (вблизи торцов) эффект текстурирования ослабевает вследствие малого уровня фонового давления.

\subsection{Сравнительный анализ интегральных показателей}

Интегральные характеристики гладкого подшипника и подшипника с суперэллиптическими углублениями для выбранных значений эксцентриситета представлены в таблицах~\ref{tab:smooth_results} и~\ref{tab:textured_results}.

\begin{table}[H]
\small
\centering
\caption{Характеристики гладкого подшипника}
\label{tab:smooth_results}
\begin{tabular}{|c|c|c|c|}
\hline
$\boldsymbol{\varepsilon}$ & $\boldsymbol{F}$\textbf{, Н} & $\boldsymbol{\mu}$ & $\boldsymbol{Q}$\textbf{, л/с} \\
\hline
0{,}05 & 551{,}5 & 0{,}2245 & 0{,}6852 \\
\hline
0{,}16 & 1775{,}8 & 0{,}0707 & 0{,}6843 \\
\hline
0{,}26 & 3141{,}8 & 0{,}0412 & 0{,}6825 \\
\hline
0{,}37 & 4805{,}2 & 0{,}0282 & 0{,}6801 \\
\hline
0{,}48 & 6984{,}6 & 0{,}0207 & 0{,}6772 \\
\hline
0{,}59 & 10124{,}4 & 0{,}0157 & 0{,}6744 \\
\hline
0{,}69 & 15347{,}9 & 0{,}0118 & 0{,}6717 \\
\hline
0{,}80 & 25964{,}3 & 0{,}0085 & 0{,}6698 \\
\hline
\end{tabular}
\end{table}

\begin{table}[H]
\small
\centering
\caption{Характеристики подшипника с \varDimpleTypeGenitive{} углублениями}
\label{tab:textured_results}
\begin{tabular}{|c|c|c|c|}
\hline
$\boldsymbol{\varepsilon}$ & $\boldsymbol{F}$\textbf{, Н} & $\boldsymbol{\mu}$ & $\boldsymbol{Q}$\textbf{, л/с} \\
\hline
0{,}05 & 653{,}3 & 0{,}1884 & 0{,}6899 \\
\hline
0{,}16 & 1894{,}9 & 0{,}0658 & 0{,}6890 \\
\hline
0{,}26 & 3309{,}0 & 0{,}0388 & 0{,}6872 \\
\hline
0{,}37 & 5042{,}1 & 0{,}0266 & 0{,}6847 \\
\hline
0{,}48 & 7359{,}6 & 0{,}0195 & 0{,}6818 \\
\hline
0{,}59 & 10871{,}9 & 0{,}0144 & 0{,}6788 \\
\hline
0{,}69 & 17103{,}4 & 0{,}0105 & 0{,}6760 \\
\hline
0{,}80 & 32116{,}7 & 0{,}0068 & 0{,}6737 \\
\hline
\end{tabular}
\end{table}

Нагрузочная способность $F$ нелинейно возрастает с увеличением эксцентриситета для обоих вариантов подшипника. Текстурированный подшипник демонстрирует более высокую нагрузочную способность при всех значениях $\varepsilon$. При больших эксцентриситетах ($\varepsilon = 0{,}80$) эффект текстурирования проявляется особенно заметно: нагрузочная способность текстурированного подшипника (32116{,}7~Н) превышает гладкий аналог (25964{,}3~Н) на 23{,}7\%.

Коэффициент трения $\mu$ монотонно убывает с ростом $\varepsilon$ и для текстурированного подшипника оказывается ниже во всём диапазоне эксцентриситетов. Расход смазки $Q$ слабо зависит от $\varepsilon$ и у текстурированного подшипника несколько выше вследствие увеличения эффективного объёма смазочного слоя в зоне углублений.

Для количественной оценки эффективности текстурирования в таблице~\ref{tab:comparison} приведены результаты сравнения при $\varepsilon = 0{,}6$.

\begin{table}[H]
\small
\centering
\caption{Сравнение характеристик при $\varepsilon = 0{,}6$}
\label{tab:comparison}
\begin{tabular}{|l|c|c|c|}
\hline
\textbf{Характеристика} & \textbf{Гладкий} & \textbf{С углублениями} & $\boldsymbol{\Delta}$\textbf{, \%} \\
\hline
Нагрузка $F$, Н & 10685{,}6 & 11464{,}0 & $+7{,}28$ \\
\hline
Коэфф. трения $\mu$ & 0{,}01507 & 0{,}01389 & $-7{,}80$ \\
\hline
Расход $Q$, л/с & 0{,}6739 & 0{,}6785 & $+0{,}67$ \\
\hline
\end{tabular}
\end{table}

Нанесение суперэллиптических углублений повышает несущую способность подшипника на 7{,}28\% благодаря микрогидродинамическому эффекту. Коэффициент трения снижается на 7{,}80\% за счёт увеличения средней толщины смазочного слоя в области углублений. Расход смазки возрастает на 0{,}67\%, что способствует улучшению теплоотвода без значительного увеличения требуемой производительности системы маслоснабжения.

На рисунках~\ref{fig:load_vs_eps}--\ref{fig:flow_vs_eps} представлены зависимости интегральных характеристик от эксцентриситета для обоих вариантов подшипника.

\begin{figure}[H]
\centering
\includegraphics[width=\textwidth]{figures/load_vs_epsilon.png}
\caption{Нагрузочная способность $F(\varepsilon)$ для гладкого и текстурированного подшипников}
\label{fig:load_vs_eps}
\end{figure}

Кривая $F(\varepsilon)$ для текстурированного подшипника располагается выше во всём диапазоне эксцентриситетов. С ростом $\varepsilon$ разрыв между кривыми увеличивается: при $\varepsilon = 0{,}05$ разница составляет около 100~Н, тогда как при $\varepsilon = 0{,}80$ она достигает более 6000~Н. Нелинейный характер зависимости обусловлен степенной зависимостью давления от толщины зазора ($P \sim h^{-3}$), что приводит к всё более выраженному эффекту текстурирования при малых зазорах.

\begin{figure}[H]
\centering
\includegraphics[width=\textwidth]{figures/friction_vs_epsilon.png}
\caption{Коэффициент трения $\mu(\varepsilon)$ для гладкого и текстурированного подшипников}
\label{fig:friction_vs_eps}
\end{figure}

Коэффициент трения текстурированного подшипника ниже гладкого при всех значениях $\varepsilon$. Обе кривые монотонно убывают, асимптотически приближаясь к нулю при $\varepsilon \to 1$. Наибольший абсолютный выигрыш наблюдается при малых эксцентриситетах: при $\varepsilon = 0{,}05$ разница составляет $\Delta\mu = 0{,}036$, тогда как при $\varepsilon = 0{,}80$~--- $\Delta\mu = 0{,}002$. В относительном выражении снижение трения остаётся сопоставимым во всём диапазоне.

\begin{figure}[H]
\centering
\includegraphics[width=\textwidth]{figures/flow_vs_epsilon.png}
\caption{Расход смазки $Q(\varepsilon)$ для гладкого и текстурированного подшипников}
\label{fig:flow_vs_eps}
\end{figure}

Расход смазки текстурированного подшипника незначительно превышает расход гладкого. Разница составляет порядка 0{,}5--0{,}7\% и практически не зависит от эксцентриситета.

Таким образом, текстурирование суперэллиптическими углублениями обеспечивает одновременное повышение нагрузочной способности и снижение коэффициента трения при незначительном росте расхода смазки. Эффект усиливается при увеличении эксцентриситета, что особенно важно для тяжелонагруженных подшипников.

\subsection*{Выводы по главе 3}
\phantomsection
\addcontentsline{toc}{subsection}{Выводы по главе 3}

Выполнен расчёт характеристик радиального подшипника скольжения с суперэллиптическим микрорельефом для 15 значений эксцентриситета. Получены поля давления и зазора, построены зависимости нагрузочной способности, коэффициента трения и расхода смазки от эксцентриситета. Сравнительный анализ при $\varepsilon = 0{,}6$ выявил увеличение нагрузочной способности на $7{,}28$\%, снижение коэффициента трения на $7{,}80$\% и рост расхода на $0{,}67$\%, что подтверждает эффективность суперэллиптического текстурирования.
